\documentclass[12pt]{article}
\usepackage[utf8]{inputenc}
\usepackage[T1]{fontenc}
\usepackage{hyperref}
\title{Maiwa Language Resources for K-12 Curriculum Development}
\author{Compiled by Claude Sonnet 5.5 (claude-sonnet-5-5) for Andr\'as Kornai}
\date{}
\begin{document}
\maketitle

\section{Introduction}
Maiwa (ISO 639-3 wmm; Glottocode maiw1250) is an Austronesian language of South Sulawesi, Indonesia, in the Northern South Sulawesi group. It is closely related to Duri, Enrekang and Malimpung, but distinct from Enrekang (ptt). It is spoken mainly in the Maiwa and Bungin districts of Enrekang Regency. Speaker estimates conflict: Wikipedia gives about 50,000 (1990), while Joshua Project lists 7,000 for the Maiwa-Bungin people group in Indonesia and 11,000 on its language page. It is not known to be taught in schools. The search results showed little beyond encyclopedic and catalogue entries.

\section{Primary Educational Resources}
None were found in search results.

\section{Linguistic and Reference Materials}
\begin{itemize}
\item Wikipedia, Maiwa language (Sulawesi), a stub: \url{https://en.wikipedia.org/wiki/Maiwa_language_(Sulawesi)}
\item Glottolog, Maiwa: \url{https://glottolog.org/resource/languoid/id/maiw1250.bigmap.html}
\item Wikipedia, Enrekang Regency: \url{https://en.wikipedia.org/wiki/Enrekang_Regency}
\item Wikipedia, Enrekang language (a related but separate language): \url{https://en.wikipedia.org/wiki/Enrekang_language}
\end{itemize}

\section{Religious and Community Resources}
Joshua Project reports that scripture is limited to Bible portions and that the language is written. It lists Bungin and Maroangin as dialects. I saw no URL for the scripture texts themselves.
\begin{itemize}
\item Maiwa-Bungin profile: \url{https://joshuaproject.net/people_groups/20325}
\item Language page: \url{https://joshuaproject.net/languages/wmm}
\end{itemize}

\section{Digital and Audio Resources}
None were found in search results.

\section{Curriculum Development Context}
All adults speak Maiwa but not all young people do, and it is not known to be taught in schools. No orthography, primer or school materials were found. Curriculum development would have to start from community consultation and fieldwork. Regional Indonesian-language educational resources and the related languages above may help as comparative material.

\end{document}